\clearpage\fancyhead[L]{\small KOMET.}\SourceLeaf{387}
\begin{center}\textbf{KOMET.}\par\vspace{4pt}\hrule\vspace{4pt}\fontsize{10}{12}\selectfont Göttingische gelehrte Anzeigen. Stück 45. Seite 441. 20. März 1815.\par\vspace{4pt}\hrule\vspace{4pt}\end{center}

\par\indent In einem Schreiben vom 7. März an Professor \textsc{Gauss} berichtete Dr. \textsc{Olbers}, dass er am 6. März im Perseus einen neuen Kometen entdeckt hatte. Am 6. März um $10^{\mathrm u}\,55^{\mathrm m}$ mittlerer Zeit betrug seine Rektaszension $49^\circ\,7'$ und seine nördliche Deklination $32^\circ\,7'$. Am 7. März um $7^{\mathrm u}\,40^{\mathrm m}$ betrug seine Rektaszension $49^\circ\,22'$ und seine nördliche Deklination $32^\circ\,32'$. Danach ist zu erwarten, dass dieser Komet noch ziemlich lange sichtbar bleiben wird. Im Übrigen ist er sehr klein, hat einen schlecht abgegrenzten Kern und einen sehr blassen, durchsichtigen Nebel. In Helligkeit und Form ähnelt er dem zweiten Kometen von 1811, wie dieser im Februar 1812 erschien.

\par\indent Gleich nach Eingang dieses Briefs am 13. März suchte man an der hiesigen Sternwarte nach dem neuen Kometen. Nur für wenige Minuten klarte der Himmel so weit auf, dass der Komet zwischen 159 und 164 im Perseus zu sehen war (nach \textsc{Bodes} Katalog). Doch die Regenwolken kamen zu schnell zurück, sodass nicht einmal eine Schätzung möglich war.
